\GalSourcePageStart{080}{L0079}{51}{51}


\begin{center}
{\scriptsize\bfseries DES ÉQUATIONS PRIMITIVES\par}
\vspace{0.45em}
{\scriptsize\bfseries QUI SONT SOLUBLES PAR RADICAUX \GalHistoricalFootnoteMark{*}\par}
\vspace{0.45em}
{\scriptsize\bfseries (Fragment.)\par}
\end{center}
\GalHistoricalFootnoteText{*}{\emph{Voir} la Note 1 de la page 33.}

Cherchons, en général, dans quel cas une équation primitive
est soluble par radicaux. Or, nous pouvons de suite établir un
caractère général fondé sur le degré même de ces équations. Ce
caractère est celui-ci: {\itshape Pour qu'une équation primitive soit résoluble
par radicaux, il faut que son degré soit de la forme~$p^{\nu}$, $p$~étant
premier.} Et de là suivra immédiatement que, lorsqu'on aura
à résoudre par radicaux une équation irréductible dont le degré
admettrait des facteurs premiers inégaux, on ne pourra le faire
que par la méthode de décomposition due à M.~Gauss; sinon
l'équation sera insoluble.

Pour établir la propriété générale que nous venons d'énoncer
relativement aux équations primitives qu'on peut résoudre par
radicaux, nous pouvons supposer que l'équation que l'on veut
résoudre soit primitive, mais cesse de l'être par l'adjonction d'un
simple radical. En d'autres termes, nous pouvons supposer que,
$n$ étant premier, le groupe de l'équation se partage en $n$ groupes
irréductibles conjugués, mais non primitifs. Car, à moins que le
degré de l'équation soit premier, un pareil groupe se présentera
toujours dans la suite des décompositions.

Soit $N$ le degré de l'équation, et supposons qu'après une extraction
de racine de degré premier~$n$, elle devienne non primitive et
se partage en $Q$~équations primitives de degré~$P$, au moyen d'une
seule équation de degré~$Q$.

Si nous appelons $G$ le groupe de l'équation, ce groupe devra se
partager en $n$~groupes conjugués non primitifs, dans lesquels les
lettres se rangeront en systèmes composés de $P$~lettres conjointes
chacun. Voyons de combien de manières cela pourra se faire.

Soit $H$ l'un des groupes conjugués non primitifs. Il est aisé de
